

We situate our contribution across three areas of prior work.

\subsection{White-Box Probes and Black-Box Behavioral Detection}

Deception-related information is encoded in model activations and extractable with weight access---honesty/deception directions~\cite{zou2023representation}, unsupervised truthfulness probes~\cite{burns2022discovering}, activation classifiers (SAPLMA;~\citealp{azaria2023internal}), inference-time steering (ITI;~\citealp{li2024inference})---all requiring access unavailable for API-only models. Unlike externally induced misbehavior~\cite{goodfellow2014explaining, perez2022ignore}, the model itself is the source of the deceptive output here.

\paragraph{Black-box behavioral approaches.} The closest prior work is \citet{pacchiardi2023catch}, who ask a fixed battery of \emph{unrelated} follow-up questions after a model states a possibly false claim and train a classifier on the binarised answers, finding detectable changes even on topically unrelated queries. We differ in three ways: (1)~our own instrument uses topically related, adaptive follow-ups; (2)~we add prompt-equalized and cross-family controls, under which most of the above-chance accuracy reported in this setting does not survive; and (3)~we give the identifiability analysis their fixed-condition evaluation leaves open, rather than a further detector. Their unrelated-question design leaves refusal-count accuracy nearly unchanged, consistent with the signal residing in the target's \emph{first} response (EXP-K, Appendix~\ref{app:pacchiardi}).

\paragraph{Instruction-following and judge controls.} \citet{sharma2023sycophancy} show LLMs shift answers toward perceived user preferences, controlling for it by comparing behavior with and without opinion-carrying context---the confound our prompt-equalized design targets, and we find the same pattern (93.9--100\% instructed $\to$ 52--69\% equalized under cross-family extraction). LLM evaluators also carry systematic biases, which motivates the cross-family control: \citet{verga2024replacing} find diverse judge panels more reliable than single judges, and here extractor choice moves accuracy by 9--16\,pp (EXP-H).

\subsection{Red-Teaming, Evaluation, and Autonomous Deception}

Red-teaming, human~\cite{ganguli2022red} and automated~\cite{perez2022red}, focuses on \emph{eliciting} harmful outputs rather than \emph{detecting} whether ostensibly normal outputs are deceptive---the complementary problem studied here.

\paragraph{Autonomous deception.} Deceptive behavior also arises without explicit instruction: strategic deception under pressure followed by misreporting~\cite{scheurer2024strategic}, alignment faking~\cite{greenblatt2024alignment}, learned reward tampering~\cite{denison2024sycophancy}, and capabilities that emerge with scale~\cite{hagendorff2024deception}. This regime matters most for safety and is the one our setting does not cover---it is also intervention~(A) of §\ref{sec:identification}, the manipulation that would identify the effect: a model deceiving autonomously need not produce the instruction-following behavior that instructed benchmarks cannot separate from deception.

\paragraph{Concurrent evaluations of lie detectors.} A recent line of work stress-tests lie detectors across settings rather than auditing one. \citet{kretschmar2025liarsbench} assemble 72{,}863 lie and honest-response examples over seven datasets and find black- and white-box techniques fail systematically on lie types not adjudicable from the transcript alone. \citet{cooney2026didyoulie} evaluate four detectors on 31 models (2B--1T) and 13 \emph{trained} organisms with verified hidden beliefs: all four scale positively on \emph{prompted} lying, while activation- and logprob-based detectors drop sharply on the organisms. \citet{hopkins2026liedetectors} report the same for fine-tuned detectors, at AUROC 0.95 in-distribution but 0.70--0.75 across held-out lie categories, barely beating prompted baselines. Those trained organisms instance intervention~(A)---deception varying with elicitation fixed---and the prompted-to-trained gap all three report is what the Proposition predicts: accuracy earned by manipulating the instruction carries no commitment about deception arising another way.

\paragraph{Summary of the gap.} These establish that lie detectors fail to \emph{generalize} across lie types, organisms, and training regimes. None isolates the identifiability problem created when the intervention that induces deception also induces instruction-following---so none answers whether the accuracy reported in the standard paradigm was ever attributable to deception. That is the question addressed here. The constraint itself is a standard non-identifiability result; what is new is enforcing it on a widely used paradigm and turning it into the five-criterion protocol of §\ref{sec:eval_controls}.
